% Body of the example paper. Each file in venues/ sets up a document class
% and its front matter, defines \venue, and inputs this file. All numbers are
% made up for illustration.

\emergencystretch=1.5em

\makeatletter
\newenvironment{widetable}
  {\if@twocolumn\begin{table*}[t]\else\begin{table}[t]\fi\centering\small
   \setlength{\tabcolsep}{5pt}}
  {\if@twocolumn\end{table*}\else\end{table}\fi}
\newenvironment{widefigure}
  {\if@twocolumn\begin{figure*}[t]\else\begin{figure}[t]\fi\centering}
  {\if@twocolumn\end{figure*}\else\end{figure}\fi}
\makeatother
\newcommand\examplefig[1]{%
  \IfFileExists{../build/figs/\venue/#1.pdf}
    {\includegraphics[width=\linewidth]{../build/figs/\venue/#1.pdf}}
    {\fbox{\parbox{0.9\linewidth}{\centering\footnotesize
      Run \texttt{make examples} to draw this figure.}}}}

\section{Tables}

Each table below is an ordinary table body inside \texttt{pvtable}. The
package finds the numbers, marks the best and second-best value in each
column or row, and colors the marks from one accent color. The key under
each table shows every mark the table uses, so the captions only state the
finding and the setup.
Table~\ref{tab:leaderboard} compares methods across many benchmarks,
Table~\ref{tab:main} adds the change against a baseline,
Table~\ref{tab:ablation} shows only that change, and Table~\ref{tab:grid}
shades a grid of values.

\begin{widetable}
  \caption{\textbf{Ours leads on most benchmarks, but not on latency.}
    Benchmarks are rows and methods are columns.}
  \label{tab:leaderboard}
  \begin{pvtable}[by=row]{l*{5}{c}}
    \toprule
     & \pvours Ours & Alpha & Beta & Gamma & Delta \\
    \midrule
    \pvsection{Overall} \\
    Index v5 & 20.1 & 4.4 & 6.3 & 20.1 & 13.6 \\
    Index v2 & 46.5\% & -- & -- & -- & 29.8\% \\
    \pvsection{Coding} \\
    \pvlabel{CodeBench v2}{agent harness v0.2} & 15.1\% & 0.0\% & 1.5\% & 14.9\% & 6.1\% \\
    \pvlabel{RepoFix (Hard)}{agent harness v2.4} & 77.9\% & 14.2\% & 44.3\% & 73.8\% & 44.1\% \\
    \pvlabel{RepoFix (Multilingual)}{bash only} & 73.7\% & 43.0\% & 66.9\% & 70.0\% & 67.7\% \\
    \pvsection{Search and tool use} \\
    WebQA & 79.0\% & 37.0\% & 63.9\% & 61.6\% & 78.9\% \\
    ToolBench & 82.7\% & 54.2\% & 75.5\% & 77.3\% & 67.7\% \\
    Latency (s) \pvdown & 1.42 & 0.91 & 1.10 & 1.38 & 2.05 \\
    \bottomrule
  \end{pvtable}
  \pvkey
\end{widetable}

\begin{table}[t]
  \centering\small
  \caption{\textbf{Ours improves every metric over the base model.}
    Accuracy is the mean of three seeds with its standard deviation.}
  \label{tab:main}
  \begin{pvtable}[delta=absolute]{lccc}
    \toprule
    Method & Acc \pvup & WER \pvdown & RTF \pvdown \\
    \midrule
    \pvbaseline Base & 70.1 & 12.4 & 0.31 \\
    Prior A & 72.3 & 11.9 & 0.35 \\
    Prior B & 71.6 & 12.8 & 0.22 \\
    \pvours Ours & $74.0 \pm 0.4$ & 10.2 & 0.28 \\
    \bottomrule
  \end{pvtable}
  \pvkey[baseline=Base]
\end{table}

\begin{table}[t]
  \centering\small
  \caption{\textbf{Removing the adapter costs the most accuracy.}
    The full model row gives the absolute values.}
  \label{tab:ablation}
  \begin{pvtable}[delta=absolute, show=delta, ranks=0]{lcc}
    \toprule
    Variant & Acc \pvup & WER \pvdown \\
    \midrule
    \pvbaseline Full model & 74.0 & 10.2 \\
    without adapter & 70.8 & 11.6 \\
    without warm-up & 73.1 & 10.4 \\
    without data mixing & 74.2 & 10.9 \\
    \bottomrule
  \end{pvtable}
  \pvkey[baseline=the full model]
\end{table}

\begin{table}[t]
  \centering\small
  \caption{\textbf{Accuracy grows with both model size and data.}}
  \label{tab:grid}
  \begin{pvtable}[heat=table, ranks=0]{lcccc}
    \toprule
    & \multicolumn{4}{c}{Training tokens} \\
    \cmidrule(lr){2-5}
    Size & 10B & 30B & 100B & 300B \\
    \midrule
    125M & 41.2 & 44.0 & 45.1 & 45.6 \\
    350M & 45.8 & 49.3 & 51.7 & 52.4 \\
    1.3B & 48.1 & 53.9 & 57.6 & 59.8 \\
    \bottomrule
  \end{pvtable}
  \pvkey
\end{table}

\section{Figures}

Each figure is drawn by its own short script, one per kind of figure, from
the layout file this document writes. So each one is built at this venue's
column width with text sized from this venue's body font, and with the same
accent color as the tables. The scripts are in the \texttt{figures} folder
next to this paper's source.

\begin{widefigure}
  \examplefig{bars}
  \caption{\textbf{Ours scores highest on four of five benchmarks.}
    Illustrative numbers.}
  \label{fig:bars}
\end{widefigure}

\begin{figure}[t]
  \centering
  \examplefig{deltas}
  \caption{\textbf{The adapter matters most; data mixing does not clearly help.}
    Accuracy change against the full model with 95\% intervals of the
    change over five seeds.}
  \label{fig:deltas}
\end{figure}

\begin{figure}[t]
  \centering
  \examplefig{curves}
  \caption{\textbf{Ours reaches the final reward of the best baseline in
    less than half the steps.} Mean of three seeds with 95\% intervals.}
  \label{fig:curves}
\end{figure}

\begin{figure}[t]
  \centering
  \examplefig{tradeoff}
  \caption{\textbf{Ours moves the accuracy--latency frontier.} The dashed
    line joins methods that no other method beats on both measures.}
  \label{fig:tradeoff}
\end{figure}

\begin{figure}[t]
  \centering
  \examplefig{heatmap}
  \caption{\textbf{Larger models gain more from more data.} The same grid as
    Table~\ref{tab:grid}; the best value of each column is bold.}
  \label{fig:heatmap}
\end{figure}

\section{Diagrams}

Diagrams use the same colors and print modes as the tables, and are drawn at
the printed size: \texttt{pvdiagram} warns when a diagram is wider than the
column instead of letting it be scaled down. Figure~\ref{fig:transformer}
fills its blocks by kind; filling by role would leave every block white
except the part a paper adds.

\input{../diagrams/transformer}
\begin{figure}[t]
  \centering
  \transformer
  \caption{\textbf{The model is an encoder and a decoder, each a stack of
    $N$ identical layers.} Redrawn from Vaswani et al.\ (2017).}
  \label{fig:transformer}
\end{figure}
